\documentclass[10pt,a4paper]{amsart}
\usepackage[margin=19mm]{geometry}
\usepackage{amsmath,amssymb,mathtools}
\usepackage[scr=rsfs,cal=euler]{mathalfa}
\usepackage{xfrac}
\usepackage[unicode,hidelinks,bookmarks=false]{hyperref}
\newcommand{\ie}{i.e.\ }
\newcommand{\cf}{cf.\ }
\newcommand{\resp}{respectively\ }
\newcommand{\Subsection}{\subsection}
\providecommand{\nobreakdash}{\nobreak\hbox{-}}
\setlength{\emergencystretch}{2em}
\multlinegap0pt
\usepackage{amssymb}
\usepackage{tikz}
\usepackage{float}
\newtheorem{proposition}{Proposition}
\newtheorem{lemma}{Lemma}
\title{At Most Two Infinite Blue Clusters in the CMR Representation of the Edwards--Anderson Spin Glass}
\author{Yan Ru Pei}
\date{Source: arXiv:2605.17338v2 (CC BY 4.0)}
\begin{document}
\maketitle

\noindent\emph{Source and licence.} At Most Two Infinite Blue Clusters in the CMR Representation of the Edwards--Anderson Spin Glass. Version arXiv:2605.17338v2 (CC BY 4.0), distributed by arXiv under the Creative Commons Attribution 4.0 International licence, http://creativecommons.org/licenses/by/4.0/, as recorded in the arXiv licence metadata for this exact version and shown on the abstract page https://arxiv.org/abs/2605.17338v2. Full licence notice: \texttt{licenses/arxiv-2605.17338-CC-BY-4.0.txt}.\par\smallskip

\noindent\textit{Editorial scope.} Counted unit: the article's lemma lem:resample (source0.tex lines 732--754). The article's complete original proof of it is delivered with it (source0.tex lines 756--795, one \texttt{proof} environment of 40 lines). No mathematics is written, repaired or re-derived: the statement and the proof are the article's own lines, in the article's own notation, and no retained line is substituted or re-flowed.  Retained uncounted as the source-local context the proof needs: lines 313--322.  Nothing was omitted from any retained span, so the package carries no omission record.  No retained line contains a citation, so the wrapper carries no bibliography and no bibliography entry.  No retained line contains a figure environment, an image-inclusion command or drawing code, so no image is delivered and the package declares no asset dependency.  Editorial formatting only, in addition: an AMS \texttt{amsart} preamble that restates the article's own declarations and macros used by the retained lines, namely the environments \texttt{proposition}, \texttt{lemma} on separate counters, together with the standard packages those lines need.  The retained environments print in this document as Proposition 1, Lemma 1 in order of appearance, whereas the article prints the same items with the numbering of its own full document.  The complete native source of the article as distributed in the arXiv e-print for this exact version is bundled unchanged as \texttt{sources/200788/source0.tex}.
\begin{proposition}[Local equivalence of coupling conditionals]\label{prop:bayes}
Let $S \subset E$ be finite.
Under $\mathbf{P}_\beta$, the regular conditional law of $J_S$ given $\mathcal{T}_S'$ is absolutely continuous with respect to $\bigotimes_{e \in S} \mathscr{P}$, and its density is $\mathbf{P}_\beta$-a.s.\ strictly positive.
Consequently, for every measurable $A \subset \mathrm{supp}(\mathscr{P})^{|S|}$ with $\bigotimes_{e \in S} \mathscr{P}(A) > 0$,
\begin{equation}\label{eq:bayes-positive}
  \mathbf{P}_\beta\bigl(J_S \in A \mid \mathcal{T}_S'\bigr)
  \;>\; 0.
\end{equation}
In particular, by~\textup{(D2)} couplings of either sign can be sampled with positive probability.
\end{proposition}

\begin{lemma}[CMR box filling]\label{lem:resample}
Fix $s\in\{\pm1\}$.
Let $\Lambda\subset\mathbb{Z}^d$ be a finite box, and let
\[
  E_\Lambda:=\{e\in E:e\cap\Lambda\neq\emptyset\}
\]
be the set of edges touching $\Lambda$, including its boundary edges.
Let $\mathscr{F}_{\Lambda^c}$ be the sigma-field generated by the couplings and bond variables on edges with both endpoints in $\Lambda^c$ and by the spins on $\Lambda^c$.
Let $B_s(\Lambda)$ be the event that:
\begin{enumerate}
\item[\textup{(i)}] every vertex of $\Lambda$ has overlap $q_i=s$;
\item[\textup{(ii)}] every edge with both endpoints in $\Lambda$ is blue;
\item[\textup{(iii)}] every boundary edge whose outside endpoint has overlap $s$ is blue;
\item[\textup{(iv)}] every boundary edge whose outside endpoint has overlap $-s$ is non-blue.
\end{enumerate}
Then
\[
  \mathbf{P}_\beta\bigl(B_s(\Lambda)
       \mid\mathscr{F}_{\Lambda^c}\bigr)>0
  \qquad \mathbf{P}_\beta\text{-a.s.}
\]
On $B_s(\Lambda)$, all exterior blue components in the overlap class $q=s$ that meet $\partial^+\Lambda$ are joined through $\Lambda$.
\end{lemma}

\begin{proof}
Fix the outside configuration in $\mathscr{F}_{\Lambda^c}$.
We build the desired event in three steps.

\medskip\noindent
\textbf{Step 1: prescribe the spins in $\Lambda$.}
Let
\[
  (\sigma_i,\tau_i)=
  \begin{cases}
    (+1,+1),&s=+1,\\
    (+1,-1),&s=-1,
  \end{cases}
  \qquad i\in\Lambda.
\]
Conditional on all couplings, the finite-box two-replica DLR kernel assigns strictly positive weight to every spin pattern.
Averaging over the unfrozen couplings on $E_\Lambda$ preserves strict positivity, so this event has positive conditional probability given $\mathscr{F}_{\Lambda^c}$.

\medskip\noindent
\textbf{Step 2: prescribe the couplings on $E_\Lambda$.}
On the event from Step~1, require
\begin{itemize}
\item $J_e>0$ for every edge $e$ with both endpoints in $\Lambda$;
\item $J_e\sigma_j>0$ for every boundary edge $e=\{i,j\}$ with $i\in\Lambda$ and $j\in\Lambda^c$.
\end{itemize}
This sign pattern has positive $\mathscr{P}^{\otimes E_\Lambda}$-measure by~\textup{(D2)}.
Once the spins and exterior configuration have been frozen, the conditioning field is exactly $\mathcal{T}'_{E_\Lambda}$ from Proposition~\ref{prop:bayes}.
That proposition therefore gives the prescribed sign pattern positive conditional probability.

\medskip\noindent
\textbf{Step 3: realize the bond states on $E_\Lambda$.}
Given the preceding spin and coupling choices, every interior edge is satisfied by both replicas.
A boundary edge is satisfied by both replicas exactly when its outside endpoint has overlap $s$; when that endpoint has overlap $-s$, exactly one replica satisfies it.
Thus every edge required to be blue in~\textup{(ii)}--\textup{(iii)} is blue with probability
$1-e^{-4\beta|J_e|}>0$, while the edges in~\textup{(iv)} are automatically non-blue.
The bond variables are conditionally independent given $(J,\sigma,\tau)$, and $E_\Lambda$ is finite, so the conjunction has positive conditional probability.

Successive conditioning through the three steps proves the displayed claim.
The joining statement follows from~\textup{(i)}--\textup{(iii)}.
\end{proof}

\end{document}
